\chapter{Direct Outreach to Startup Founders}
\label{ch:direct-outreach}

This is the most tactical chapter in the guide: instead of waiting for vacancies to appear, you identify funded
startups and contact the people who can influence hiring decisions directly. It operationalizes the $N$ (number
reached), $A$ (acceptance), $R$ (reply), $P$ (positive reply), $C$ (call), and $V$ (conversion) variables from
Chapter~\ref{ch:outreach-engineering}.

\section{Step 1 --- Identify Relevant Funded Startups}

Look for companies that have recently raised funding and have a product or business problem matching your
skills. Favor startups with small teams, where founders are potentially accessible and hiring decisions can be
flexible. Practical selection criteria: the company has a relevant technical or business need; your skills align
with its product or infrastructure; you can identify a useful improvement or feature; you can actually reach a
founder or other relevant decision-maker; the company's business and hiring arrangements are compatible with your
location and work authorization.

\begin{caution}
There is no single verified database that guarantees this list for you. Build it deliberately --- funding
announcements (Crunchbase, TechCrunch, sector newsletters), LinkedIn's own company and ``funding'' filters, and
niche job boards for funded startups are reasonable starting points. If you have access to a pre-built leads list
(for example, a spreadsheet of recently funded companies in a specific market), treat it as a starting point to
filter by \emph{relevance} to your skills, not as a list to spray messages at indiscriminately --- untargeted
volume is exactly what collapses the $A$ and $R$ variables toward zero in the outreach equation.
\end{caution}

\section{Step 2 --- Research the Company Before Contacting It}

Find a specific, genuine reason to contact the founder --- a recently launched support product, an AI-powered
workflow they run, or an identifiable technical bottleneck. Connect the outreach message to a real business
problem rather than sending a generic request for employment.

\section{Step 3 --- Provide Value Before Asking for a Job}

This is the main differentiator in the method.

\begin{example}
\bad \textbf{Poor message:} ``Hi, I am looking for opportunities in ML engineering at your company. Can we
connect?'' --- ineffective because it describes what you want without giving the founder a reason to respond.
\end{example}

The alternative: identify something useful you could build or improve --- a prototype for routing customer-support
tickets, a demonstration of an AI feature relevant to their product, a small workflow improvement, or an analysis
of a technical problem relevant to their business. The message should make the connection between your work and
their needs immediately obvious.

\section{Step 4 --- Use a Small, Specific Call to Action}

Offer a prototype, or propose to build one within roughly 48 hours. Make the next step easy for the founder to
evaluate --- ask whether a small demonstration would be useful, rather than asking for a formal interview or a
long meeting. This reduces the effort required to assess your work, though it does not guarantee a response.

\begin{template}
\textbf{Illustrative outreach template}\\[4pt]
Hi [Founder's Name],\\[2pt]
I saw that [Company] recently launched or is working on [specific product or feature]. I was exploring how
[specific workflow or problem] could be improved using [relevant technical approach].\\[2pt]
I have an idea for a small prototype that could help with [specific, plausible benefit]. I can share a working
demonstration and the technical trade-offs behind it.\\[2pt]
Would this be useful to your team?\\[2pt]
Best,\\{}
[Your Name]
\end{template}

\begin{caution}
This is an illustrative structure, not a script to copy verbatim --- the strongest version references a problem
you have genuinely researched and an improvement you can reasonably deliver. (Compare this to the ``break the
expectation'' message anatomy in \S2.4 --- both sources converge on the same point: generic asks get ignored,
specific, researched value gets replies.)
\end{caution}

\section{Why Prototypes Might Help: The Reciprocity Principle}

Delivering something useful before asking for an opportunity can make a founder more willing to engage --- when a
candidate invests time in understanding the company and creates something relevant, the founder has more evidence
to evaluate than a resume alone provides. The proposed sequence:

\begin{enumerate}
  \item \textbf{Research a company} --- identify a concrete, relevant problem.
  \item \textbf{Build a useful prototype} --- demonstrate technical capability and judgment.
  \item \textbf{Contact the founder} --- explain the problem, the solution, and the potential value.
  \item \textbf{Discuss the opportunity} --- the founder may assess the project directly or connect you with the
  relevant team.
\end{enumerate}

Publicly posting about the project and tagging the founder can help other companies in the same sector notice
similar work --- but this should be done professionally, is not essential, and a prototype must never expose
confidential information, copy proprietary functionality, or be treated as entitlement to a job. The goal is to
demonstrate initiative and capability, not to work indefinitely unpaid.

\begin{caution}
\textbf{Prototype-first outreach has a real cost.} Building an individual project for every company can consume
a huge amount of time. A more efficient pattern is usually to build one strong, reusable project (see
Chapter~\ref{ch:projects}) and \emph{adapt its demonstration} to each company, building a bespoke small prototype
only after a genuinely promising opportunity is identified.
\end{caution}

\section{Combine Multiple Outreach Channels}

Contact founders through several channels: LinkedIn messages, direct email, X (formerly Twitter) or another
relevant professional platform, and public technical posts that demonstrate your work. Repeated, varied exposure
makes you easier to recognize --- this is the $L$ (LinkedIn Authority) multiplier again, reinforcing the same
person's awareness of you through more than one surface.

\begin{caution}
This is not license to repeatedly spam everywhere. Use personalized, spaced follow-ups (every 2--3 days, per
\S2.4), respect requests not to be contacted, and avoid anything that reads as spam --- a sensible cadence beats
raw volume.
\end{caution}

\section{Sizing the Funnel: What to Expect}

Chapter~\ref{ch:outreach-engineering} already gave the mathematical model; here is the illustrative numeric funnel
from the remote-jobs playbook, for calibration against your own tracked numbers:

\begin{longtable}{p{0.3\textwidth} r p{0.34\textwidth}}
\toprule
\textbf{Stage} & \textbf{Illustrative count} & \textbf{Implied conversion rate} \\
\midrule
\endhead
Initial outreach (people contacted in a month) & 800 & --- \\
Conversations & 40 & 5\% of outreach \\
Interested prospects & 4 & 10\% of conversations \\
Successful hire (target outcome) & 1 & 25\% of interested prospects \\
\emph{Outreach to successful hire, overall} & --- & \emph{0.125\%} \\
\bottomrule
\end{longtable}

\begin{caution}
This is a hypothetical funnel from one speaker, not an established industry conversion rate. Treat job-seeking as
a repeatable process --- track outreach volume, responses, conversations, and outcomes, then improve the process
--- rather than expecting any single message to succeed. In practice, prioritize relevant, personalized outreach
over sheer volume: 800 low-quality messages can underperform a much smaller number of well-researched approaches,
since quality lifts every multiplier ($A$, $R$, $P$, $C$, $V$) simultaneously rather than just $N$.
\end{caution}

